%! Right Angles Revisited --- Setting Up Orthogonality — Unit 4, Lesson 4.0 — Group Activity (Answer Key)
\documentclass[10pt]{article}
\usepackage{linalg-article}
\usepackage{linalg-key}   % pulls in -boxes; do NOT also load -boxes
\newcommand{\TallMath}[1]{$\displaystyle #1\rule[-1.4em]{0pt}{3.2em}$}
\newcommand{\vv}{\mathbf{v}}
\newcommand{\ww}{\mathbf{w}}
\newcommand{\avec}{\mathbf{a}}
\newcommand{\bb}{\mathbf{b}}
\newcommand{\xx}{\mathbf{x}}
\newcommand{\pp}{\mathbf{p}}
\newcommand{\zero}{\mathbf{0}}

\begin{document}

\pageheader{Unit 4, Lesson 4.0}{Group Activity --- Answer Key}
\namepartnerperiod

\vspace{0.10in}

\begin{headlinebox}{blush}
\small\textbf{Right Angles, Everywhere.} Two vectors are \textbf{perpendicular} exactly when their
\textbf{dot product is $0$}. That one rule is the whole of Unit~4. Work with your partner, show each
dot product, and \emph{justify} every ``perpendicular'' claim.
\end{headlinebox}

\vspace{0.10in}

\begin{tcolorbox}[colback=white, colframe=black!40, title=\textbf{Tier R --- Remediate},
    fonttitle=\bfseries, arc=1.5mm, left=3mm, right=3mm, top=2mm, bottom=2mm, breakable]
Warm up the three recalled skills. Show your work.
\begin{enumerate}[leftmargin=1.7em, itemsep=8pt]
  \item Compute the dot product
        $\begin{bsmallmatrix} 1\\2 \end{bsmallmatrix}\cdot\begin{bsmallmatrix} 3\\1 \end{bsmallmatrix}
        = \ans{$5$}$. Is the angle acute, right, or obtuse? \; \ans{acute (positive)}
  \item Find the length and unit vector of $\begin{bsmallmatrix} 8\\6 \end{bsmallmatrix}$:
        \quad $\|\cdot\| = \ans{$10$}$, \quad
        unit vector $= \begin{bmatrix}{\color{keyred}\mathbf{4/5}}\\[2pt]{\color{keyred}\mathbf{3/5}}\end{bmatrix}$.
  \item Are $\begin{bsmallmatrix} 2\\6 \end{bsmallmatrix}$ and $\begin{bsmallmatrix} 3\\-1 \end{bsmallmatrix}$
        perpendicular? Compute the dot product and decide. \; \ans{Yes --- $6-6=0$}
\end{enumerate}
\end{tcolorbox}

\begin{tcolorbox}[colback=white, colframe=black!40, title=\textbf{Tier A --- Approaching Proficiency},
    fonttitle=\bfseries, arc=1.5mm, left=3mm, right=3mm, top=2mm, bottom=2mm, breakable]
Now \emph{solve for} and \emph{build} orthogonality.
\begin{enumerate}[leftmargin=1.7em, itemsep=9pt]
  \item \textbf{Solve for the entry.} Find $x$ so that $\begin{bsmallmatrix} 6\\x \end{bsmallmatrix}$
        is perpendicular to $\begin{bsmallmatrix} 2\\-3 \end{bsmallmatrix}$.
        \ansline{$6(2)+x(-3)=0\Rightarrow 12-3x=0\Rightarrow x=4$.}
  \item \textbf{Build one.} Give a nonzero vector perpendicular to
        $\begin{bsmallmatrix} 3\\5 \end{bsmallmatrix}$, and check your dot product is $0$.
        (Hint: swap the entries and negate one.)
        \ansline{$\begin{bsmallmatrix} 5\\-3 \end{bsmallmatrix}$ (or $\begin{bsmallmatrix} -5\\3 \end{bsmallmatrix}$): $3(5)+5(-3)=15-15=0$. \checkmark}
  \item \textbf{The §3.6 pairing.} Let $A=\begin{bmatrix} 1 & 3 \\ 2 & 6 \end{bmatrix}$. Its
        nullspace is all multiples of $\xx=\begin{bsmallmatrix} -3\\1 \end{bsmallmatrix}$ (both rows
        give $x_1+3x_2=0$). Dot \emph{each row} of $A$ with $\xx$. What do you get, and what does it
        say about the angle between the nullspace and the rows?
        \ansline{Row~1: $(1)(-3)+(3)(1)=0$. Row~2: $(2)(-3)+(6)(1)=0$. Both $0$, so $\xx$ is perpendicular to every row --- the nullspace $\perp$ the row space (the §3.6 pairing).}
\end{enumerate}
\end{tcolorbox}

\begin{tcolorbox}[colback=white, colframe=black!40, title=\textbf{Tier E --- Extension},
    fonttitle=\bfseries, arc=1.5mm, left=3mm, right=3mm, top=2mm, bottom=2mm, breakable]
\textbf{A first look at projection.} A target $\bb=\begin{bsmallmatrix} 1\\3 \end{bsmallmatrix}$ is
\emph{not} on the line through $\avec=\begin{bsmallmatrix} 3\\1 \end{bsmallmatrix}$. We want the point
$t\avec$ on the line \emph{closest} to $\bb$ --- the projection. The closest point is where the error
$\bb-t\avec$ makes a \textbf{right angle} with the line.
\begin{enumerate}[leftmargin=1.7em, itemsep=9pt]
  \item Set the error perpendicular to $\avec$: write the equation $\avec\cdot(\bb-t\avec)=0$, then expand
        it to $\avec\cdot\bb - t\,(\avec\cdot\avec)=0$.
        \ansline{$\avec\cdot\bb=(3)(1)+(1)(3)=6$ and $\avec\cdot\avec=(3)(3)+(1)(1)=10$, so $6 - 10t = 0$.}
  \item Solve for $t$. (You should get $t=\dfrac{\avec\cdot\bb}{\avec\cdot\avec}$.) Compute its value here.
        \ansline{$t=\dfrac{\avec\cdot\bb}{\avec\cdot\avec}=\dfrac{6}{10}=0.6$.}
  \item Find the projection $\pp=t\avec$ and the error $\bb-\pp$, then \emph{check} that
        $\avec\cdot(\bb-\pp)=0$. This formula is Lesson~4.2.
        \ansline{$\pp=0.6\begin{bsmallmatrix} 3\\1 \end{bsmallmatrix}=\begin{bsmallmatrix} 1.8\\0.6 \end{bsmallmatrix}$; error $\bb-\pp=\begin{bsmallmatrix} -0.8\\2.4 \end{bsmallmatrix}$; check $\avec\cdot(\bb-\pp)=3(-0.8)+1(2.4)=-2.4+2.4=0$. \checkmark}
\end{enumerate}
\end{tcolorbox}

\begin{teachernote}
Tier~A item~3 and all of Tier~E are the point of the unit: orthogonality ($\cdot=0$) is what makes
the four-subspace pairing true \emph{and} what finds the closest point. Tier~E derives the
projection formula $t=\frac{\avec\cdot\bb}{\avec\cdot\avec}$ by \emph{forcing} the error perpendicular ---
exactly the move Lesson~4.2 makes. If time is short, Tier~E item~1 (setting up
$\avec\cdot(\bb-t\avec)=0$) is the must-do; the algebra can carry over as homework.
\end{teachernote}

\end{document}
